\documentclass[12pt,a4paper]{article}

\usepackage{graphicx, subfigure}
\usepackage{appendix}

\usepackage{listings}
\lstset{ %
language=C,                % choose the language of the code
basicstyle=\footnotesize,       % the size of the fonts that are used for the code
numbers=left,                   % where to put the line-numbers
numberstyle=\footnotesize,      % the size of the fonts that are used for the line-numbers
stepnumber=1,                   % the step between two line-numbers. If it is 1 each line will be numbered
numbersep=5pt,                  % how far the line-numbers are from the code
%backgroundcolor=\color{white},  % choose the background color. You must add \usepackage{color}
showspaces=false,               % show spaces adding particular underscores
showstringspaces=false,         % underline spaces within strings
showtabs=false,                 % show tabs within strings adding particular underscores
frame=single,           % adds a frame around the code
tabsize=2,          % sets default tabsize to 2 spaces
captionpos=b,           % sets the caption-position to bottom
breaklines=true,        % sets automatic line breaking
breakatwhitespace=false,    % sets if automatic breaks should only happen at whitespace
escapeinside={\%*}{*)}          % if you want to add a comment within your code
}


\newcommand{\HRule}{\rule{\linewidth}{0.5mm}}

\begin{document}
\begin{titlepage}

\begin{center}

% Upper part of the page
\includegraphics[scale=0.4]{figures/logo}\\[1cm]

\textsc{\large Advanced Networking 2011-2012}\\[1.5cm]

\textsc{\large Final Report}\\[0.5cm]

% Title
\HRule \\[0.5cm]
{\Large \bfseries TCP-transmission extension in\\ GRAPES library}\\[0.5cm]
%--------------------------------------------------------------------------\\[0.4cm]
%{\LARGE \bfseries Trusted Organizational Base}\\
\HRule \\[1.5cm]

% Author and supervisor
\begin{minipage}{0.4\textwidth}
\begin{flushleft} \large
Arb\"er \textsc{Fama}\\
\emph{arberfama@gmail.com}
\end{flushleft}
\end{minipage}
\begin{minipage}{0.4\textwidth}
\begin{flushright} \large
Giovanni \textsc{Simoni} \\
\emph{simgidacav@gmail.com}
\end{flushright}
\end{minipage}

\vfill

% Bottom of the page
{\large \today}

\end{center}

\end{titlepage}


\begin{abstract}
In this report it is explained a module-extension for \textbf{GRAPES} library.
It consists of a practical implementation of a \textit{net-helper} module that
manages the transmittion over \textbf{TCP}. In the first part, there is an
introduction to \textit{net-helper} interface and the previous version of  
\textbf{GRAPES} library; later reported some problems introduced by
managing such communications in \textbf{TCP} and finally more details
reguarding the specific implementation.

\end{abstract}


%\tableofcontents
%\listoffigures
%\listoftables

\section{Introduction}
\label{ch:intro}
Before getting lost in many technicalities, there have been reported a small
introduction to the GRAPES library. It is a generic environment for P2P
streaming. It is designed to be usable in different environments and
situations. Peer-to-peer (\textbf{P2P}) tecnologies are becoming increasingly
popular as a way to overcome the scalability problem limitations intrinsic in
traditional applications based on client/server paradigm. In particular there
is a great interest in P2P streaming applications because they have high
demands in term of bandwidth requirements and IP-level multicasting is not
supported on the internet.
More details could be found in \cite{disi10-038}.

The \textbf{GRAPES} library provides a net-helper for managing the
communication. At the beggining, it was using the \textit{UDP} protocol on POSIX
systems. As it is known, \textit{UDP} is commonly used in internet for streaming audio and video,
whereas it is avoided in contexts where the content of the transmitted data is relevant.
The reason for implementing a \textit{TCP}-communication module that works in a context where
commonly \textit{UDP} is used, is given by the need of 'guaranteeing delivery'.
Some features for this are intrinsecally given by \textit{TCP}:
\begin{itemize}
\item \textbf{flow control}: it determines when data needs to be re-sent and stops flow of data until
  previous packes are successfully transeferred.
\item \textbf{re-transimittion}: if the data is lost or collision have occur during transmition.
\end{itemize}
Some other functionalities have been created in order to takle disadvantages that \textit{TCP} has
respect to \textit{UDP}:
\begin{itemize}
\item \textbf{dictionary}: a manager for the set of connections in order to keep them `alive'.
                          Commonly \textit{TCP} is used in 'single-shot' mode: the transition is
                          used for delivering a packet and the connection is closed once it has
                          been completely transferred. This is an waste of resources in an
                          environment where is exchanged a high rate of packets.
\item \textbf{states machines}: for both sender and receiver peers have been builded a state machine
                                that permites to them to be in consistent state.
\item \textbf{timeout}: connections are closed when a predefined timeout expires.
\end{itemize}

In this report will be explained the internals of the net-helper-tcp module.
In particular, the focus will be on three main functions that are involved:
\begin{itemize}
  \item \textit{send\_to\_peer}: executes the sending of a messages. More details about
    this could be found in next section.
  \item \textit{recv\_from\_peer}: executes the receiving of a message. More details about
    this could be found in next section.
  \item \textit{wait4data}:  continuosly scans connected peers for receiving data.
    More details about this could be found in next section.
\end{itemize}


\subsection{Connection estabilishment}

\subsubsection{net-helper init}
At first, a \textit{node} is created containg the own \textit{ip address} and \textit{port}. 

\begin{lstlisting}
typedef struct nodeID {
    unsigned refcount;
    sockaddr_t addr;
    local_t *loc;
} nodeid_t;
\end{lstlisting}
This data set contains:
\begin{itemize}
\item \textbf{reference counter} a label that accounts the duplication of the node.
\item \textbf{socket address} a data structure containing the socked address
  (for more details: \ref{apx-nh-init}). 
\item \textbf{local\_t} local data struct that will be explained below.
\end{itemize}
Than there is allocated a \textit{local\_t} data structure, as mentioned before. The \textit{nodeid\_t}
contains a local data structure that is:
\begin{lstlisting}
typedef struct {
    dict_t neighbors;
    int srvfd;
    aqueue_t aqueue;
} local_t;
\end{lstlisting}
This data set consists on:
\begin{itemize}
\item \textbf{a dictionary} (a set of nodes) of its neighbours. More details about this will be explained in next
                            section.
\item \textbf{a file descriptor} that correspond to its own file descriptor.
\item \textbf{a queue} that is a data structure corresponding to a queue of active peers. More details about this
                       will be explained in section %%%
\end{itemize}
After initializing the needed data structures (e.g.local\_t, dict\_t, timeout, ...),
the initializer returns the created data-structure \textit{nodeID}.

\subsection{Dictionary}
This could be considered the key point for management of the connections. It is managed through a hash-table
which relies in an external library, \textbf{libdacav}.\\
The related data structure is declare as:
\begin{lstlisting}
struct dict {
    dhash_t *hash;
    size_t nelems;
};
\end{lstlisting}

\subsubsection{Dictionary interface}
Here are reported the headers of the functions needed for the dictionary's management.
\begin{lstlisting}
/* Return 1 to continue, 0 to stop scanning */
typedef int (* dict_foreach_t) (void *ctx, const sockaddr_t * addr,
client_t cl);

/* Destructor must support NULL as input value */
dict_t dict_new (struct tag *);

/* Returns the inner pointer to client instance */
client_t dict_search (dict_t d, const sockaddr_t *);

void dict_foreach (dict_t d, dict_foreach_t cb, void * ctx);

void dict_del (dict_t d);
\end{lstlisting}

\subsection{A-queue}

\begin{lstlisting}
struct aqueue {
    dlist_t * queue;
    fd_set read;
    fd_set write;
    int maxfd;
    int n_active;
};
\end{lstlisting}

\subsubsection{A-queue iterface}
Here has been reported the interface of the functions related to the queue management:
\begin{lstlisting}
aqueue_t aqueue_new ();


/**                                                                                                        
 * @retval 0 on normal execution;                                                                          
 * @retval 1 if user_fds are involved in selects;                                                          
 * @retval -1 on error.                                                                                    
 */
int aqueue_scan_dict (aqueue_t q, dict_t dict, int srv_fd, int *user_fds,
                      struct timeval *maxwait);

int aqueue_empty (aqueue_t q);

void aqueue_del (aqueue_t q);

client_t aqueue_next (aqueue_t q);
\end{lstlisting}

\subsection{Timeout}

\begin{lstlisting}
struct tout {
    struct timeval last_update;
    struct timeval period;
};
\end{lstlisting}

\subsubsection{Timeout interface}
Here is reported the interface of the timeout management:
\begin{lstlisting}
tout_t tout_new (const struct timeval *timeout);

tout_t tout_copy (const tout_t t);

void tout_reset (tout_t t);

int tout_expired (tout_t t);

void tout_del (tout_t t);

struct timeval * tout_remaining (tout_t t, struct timeval *diff);

unsigned tout_timeval_to_ms (const struct timeval *tval);

unsigned tout_now_ms ();
\end{lstlisting}


%\section{GRAPES}
\label{ch:grapes}
At this point lets have a look to the content of \textbf{GRAPES} library and the module we developed.

\subsection{asd}
\label{asd}


\section{Net Helper - TCP}
A first layout of the \textit{net-helper-TCP}'s module, that consist on initialization, main data structures
and management functions, were explained in the previous section. whereas its internals, that consist on
logical management of the data transfer, are explained in the follow one.

\subsection{Send to Peer}
The sending conceptually has been split in two parts:
\begin{itemize}
  \item \textbf{a state machine} 
  \item \textbf{a procedure} that manages the act of sending
\end{itemize}


\subsubsection{Sender state machine}
The state machine itself has been divided in two parts: at \textit{sender} level and at \textit{message} level.
\begin{lstlisting}
typedef enum {
    SENDER_IDLE,
    SENDER_BUSY
} sender_state_t;

typedef enum {
    SND_HEADER,
    SND_MESSAGE,
    SND_IDLE,
} state_t;
\end{lstlisting}
This is done because the sending is considered to be acomplished in two steps:
\begin{itemize}
  \item at first the sender deliveres a message containing a headder in order to inform the receiver. This is
        foundamental for the recepient in order to allocate the needed resources and give a feedback for the
        related send request.
  \item than the body of the related message is sent.
\end{itemize}

THe sender's data structure is;
\begin{lstlisting}
struct sender {
    state_t state;
    void * buffer;
    size_t buflen;
    header_t hdr;
    size_t sent;
};
\end{lstlisting}


\subsubsection{Send to peer - procedure}
The interface of this function is:
\begin{lstlisting}
int send_to_peer(const struct nodeID *self, struct nodeID *to,
                 const uint8_t *buffer_ptr, int buffer_size)
\end{lstlisting}
where the meaning of the paramenters is:
\begin{itemize}
  \item \textbf{self}: is the reference to the caller's \textit{nodeID} data-structure.
  \item \textbf{to}: is the recepient's \textit{nodeID} data-structure.
  \item \textbf{buffer\_ptr}: is the pointer to the message's buffer.
  \item \textbf{buffer\_size}: represents the size of the buffer.
\end{itemize}
At fist, the a searching is executed through the active-neighbours list. If this search fails, that means that
the two nodes were not communicating, a new connection is estabilished. Than a scanning procedure is executed
through the \textit{aqueue}.

\subsubsection{Sender - Interface}
The functions contained in the interface are:
\begin{lstlisting}
sender_t sender_new ();

sender_state_t sender_state (sender_t s);

void sender_reset (sender_t s);

int sender_subscribe (sender_t s, const msgbuf_t *msg);

int sender_run (sender_t s, int fd);

void sender_del (sender_t s);
\end{lstlisting}

The meaning of each of each of them is:
\begin{itemize}
  \item \textbf{sender\_new}: initializes a new sender.
  \item \textbf{sender\_state}: this is needed for monitoring the state of the sender.
  \item \textbf{sender\_reset}: once the trasmittion is completed, the sender is eventually reset for
                                further communications.
  \item \textbf{sender\_subscribe}: the sender need to be subscribed
  \item \textbf{sender\_run}: the sender needs to be run once it is initialized or after it has been reset.
  \item \textbf{sender\_del}: when a sender is not needed anymore, the used resources have to be released.
\end{itemize}

\subsection{Receive from Peer}
The receiver is slightly different, but the conceptual split between the \textit{state machine} and the
\textit{procedure} has been keept.

\subsubsection{Receiver state machine}
Respect to the sender, the state machine of the receiver needs to account the message is ready or the
buffer is ready (and so empty) for a new message.
\begin{lstlisting}
typedef enum {
    RECVER_MSG_READY,
    RECVER_EMPTY,
    RECVER_BUSY
} recver_state_t;

typedef enum {
    RCV_HEADER,
    RCV_MESSAGE,
    RCV_COMPLETE,
} state_t;
\end{lstlisting}

\subsubsection{Receiver state machine}
The interface of this function is:
\begin{lstlisting}
int recv_from_peer(const struct nodeID *self, struct nodeID **remote,
                   uint8_t *buffer_ptr, int buffer_size)
\end{lstlisting}
where the meaning of the paramenters is:
\begin{itemize}
  \item \textbf{self}: is the reference to the caller's \textit{nodeID} data-structure.
  \item \textbf{to}: is the sender's \textit{nodeID} data-structure.
  \item \textbf{buffer\_ptr}: is the pointer to the message's buffer.
  \item \textbf{buffer\_size}: represents the size of the buffer.
\end{itemize}
At first, there is an \textit{while} cycle in this function that iterates meanwhile the local
\textit{qaueue} is empty:
\begin{lstlisting}
while (aqueue_empty(loc->aqueue)) {
  if (aqueue_scan_dict(loc->aqueue, loc->neighbors, loc->srvfd, NULL, NULL) == -1)
    return -1;
}
\end{lstlisting}
Than, the sender is associated to the next client in the \textit{aqueue}. After that, the mesage is extracted
from the sender's data. At this point, the receiver obtains the needed message and prepares the
related buffers. At at last, the function returns the size of the obtained message.

\subsubsection{Receiver - Interface}
\begin{lstlisting}
recver_t recver_new ();

recver_state_t recver_state (recver_t r);

void recver_reset (recver_t r);

const msgbuf_t * recver_read (recver_t r);

int recver_run (recver_t r, int fd);

void recver_del (recver_t r);
\end{lstlisting}

The meaning of each of this functions is similar to what explained in the sender's section.

\subsection{Wait for Data}
This is the point in the module where the \textit{timeout} structure, explained previously, is beeing used.
The core of this function is the waiting untill the timeout expires or there is a message:
\begin{lstlisting}
   while (!tout_expired(T) && aqueue_empty(loc->aqueue)) {
        res = aqueue_scan_dict(loc->aqueue, loc->neighbors, loc->srvfd,
                               NULL, tout_remaining(T, &tmp));
                               if (res == -1) {
            tout_del(T);
            return -1;
        }
        ret |= res;
    }
\end{lstlisting}


\section{Tests}
\label{ch:tests}
In this chapter we are going to describe the tests ....



\section{Performances Analysis}
In this chapter we are going to annalyse the performances...

\begin{center}
    \begin{tabular}{ | p{4.8cm} | p{4.8cm} | p{4.8cm} |}
    \hline
       \textbf{column 1}
         & \textbf{column 2}
         & \textbf{column 3}\\ \hline \hline
       .....
         & ------
         & .... \\ \hline
     \end{tabular}
\end{center}


\section{Conclusions}


\begin{appendices}

\section{Net-helper TCP}

\subsection{Send to peer}
\label{nh-send2peer}
\begin{lstlisting}
int send_to_peer(const struct nodeID *self, struct nodeID *to,
                 const uint8_t *buffer_ptr, int buffer_size)
{
    client_t cl;
    local_t *loc;
    struct timeval zero = {
        .tv_sec = 0,
        .tv_usec = 0
    };
    const msgbuf_t msg = {
        .data = buffer_ptr,
        .size = buffer_size
    };

    assert(self->loc != NULL);

    loc = self->loc;

    cl = dict_search(loc->neighbors, &to->addr);
    if (!client_valid(cl)) {
        int clfd;

        if (client_connect(cl, &to->addr) == -1) {
            return -1;
        }
        clfd = client_get_fd(cl);
        if (sockaddr_send_hello(&self->addr, clfd) == -1) {
            close(clfd);
            return -1;
        }
    }

    if (client_write(cl, &msg) == -1) return -1;
    if (aqueue_scan_dict(loc->aqueue, loc->neighbors,
                         loc->srvfd, NULL, &zero) == -1) {
        return -1;
    }
    return buffer_size;
}
\end{lstlisting}

\subsection{Receive from peer}
\label{nh-rcv-peer}
\begin{lstlisting}
int recv_from_peer(const struct nodeID *self, struct nodeID **remote,
                   uint8_t *buffer_ptr, int buffer_size)
{
    local_t *loc;
    client_t sender;
    nodeid_t *sender_node;
    const msgbuf_t *msg;

    assert(self->loc != NULL);
    loc = self->loc;

    while (aqueue_empty(loc->aqueue)) {
        if (aqueue_scan_dict(loc->aqueue, loc->neighbors, loc->srvfd,
                             NULL, NULL) == -1) {
            return -1;
        }
    }

    sender = aqueue_next(loc->aqueue);
    msg = client_read(sender);
    assert(msg != NULL);

    if (msg->size > buffer_size) {
        print_err(``recv_from_peer'', NULL, ENOBUFS);
        return -1;
    }
    memcpy((void *)buffer_ptr, msg->data, msg->size);

    sender_node = create_node(NULL, 0);
    sockaddr_copy(&sender_node->addr, client_get_remote(sender));
    *remote = sender_node;

    return (int) msg->size;
}
\end{lstlisting}

\subsection{Wait for data}
\label{nh-wait4data}
\begin{lstlisting}
int wait4data(const struct nodeID *self, struct timeval *tout,
              int *user_fds)
{
    local_t *loc;
    struct timeval tmp;
    int res, ret;
    tout_t T;

    assert(self->loc != NULL);
    loc = self->loc;

    T = tout_new(tout);
    ret = 0;
    while (!tout_expired(T) && aqueue_empty(loc->aqueue)) {
        res = aqueue_scan_dict(loc->aqueue, loc->neighbors, loc->srvfd,
                               NULL, tout_remaining(T, &tmp));
        if (res == -1) {
            tout_del(T);
            return -1;
        }
        ret |= res;
    }
    tout_del(T);

    return ret || !aqueue_empty(loc->aqueue);
}
\end{lstlisting}


\subsection{Other functions}
\label{nh-oth-func}
\begin{lstlisting}
int tcp_serve (int backlog, sockaddr_t *addr)
{
    int fd;

    fd = socket(AF_INET, SOCK_STREAM, 0);
    if (fd == -1) {
      print_err(``tcp_serve'', ``socket'', errno);
        return -1;
    }

    if (bind(fd, &addr->sa, sizeof(struct sockaddr_in))) {
        print_err(``tcp_serve'', ``bind'', errno);
        close(fd);
        return -1;
    }

    if (listen(fd, backlog) == -1) {
        print_err(``tcp_serve'', ``listen'', errno);
        close(fd);
        return -1;
    }

    return fd;
}
\end{lstlisting}

\begin{lstlisting}
struct nodeID *nodeid_dup (struct nodeID *s)
{
    s->refcount ++;
    return s;
}
\end{lstlisting}

\begin{lstlisting}
int nodeid_equal (const nodeid_t *s1, const nodeid_t *s2)
{
    if (s1 == s2) return 1;
    return sockaddr_equal(&s1->addr, &s2->addr);
}
\end{lstlisting}

\begin{lstlisting}
struct nodeID *create_node (const char *ipaddr, int port)
{
    nodeid_t *ret = mem_new(sizeof(struct nodeID));

    /* Initialization of address */
    memset(&ret->addr, 0, sizeof(sockaddr_t));
    ret->addr.sa.sa_family = AF_INET;
    sockaddr_set_port(&ret->addr, port);

    if (ipaddr != NULL) {
        if (sockaddr_in_init(&ret->addr, ipaddr, port) == -1) {
            return NULL;
        }
    }

    ret->loc = NULL;
    ret->refcount = 1;

    return ret;
}
\end{lstlisting}

\begin{lstlisting}
void nodeid_free (struct nodeID *s)
{
    if (s == NULL) return;

    s->refcount --;
    if (s->refcount == 0) {
        local_t *local = s->loc;

        if (local != NULL) {
            dict_del(local->neighbors);
            if (local->srvfd != -1) close(local->srvfd);
            aqueue_del(local->aqueue);
            free(local);
        }
        free(s);
    }
}
\end{lstlisting}

\begin{lstlisting}
struct nodeID * net_helper_init (const char *ipaddr, int port,
                                 const char *config)
{
    nodeid_t *self;
    struct tag *cfg_tags;
    int backlog;
    local_t *local;

    self = create_node(ipaddr, port);
    if (self == NULL) {
        return NULL;
    }

    self->loc = local = mem_new(sizeof(local_t));

    /* Default settings */
    backlog = DEFAULT_BACKLOG;

    /* Reading settings */
    cfg_tags = NULL;
    if (config) {
      cfg_tags = config_parse(config);
    }
    if (cfg_tags) {
        /* FIXME: this seems not to work! Testing needed */
        config_value_int_default(cfg_tags, CONF_KEY_BACKLOG, &backlog,
                                 DEFAULT_BACKLOG);
    }
    local->neighbors = dict_new(cfg_tags);
    free(cfg_tags);

    local->aqueue = aqueue_new();
    local->srvfd = tcp_serve(backlog, &self->addr);
    if (local->srvfd == -1) {
      nodeid_free(self);
        return NULL;
    }

    return self;
}
\end{lstlisting}

\begin{lstlisting}
struct nodeID *nodeid_undump (const uint8_t *b, int *len)
{
    nodeid_t *node;

    node = create_node(NULL, 0);
    *len = sockaddr_undump(&node->addr, sizeof(sockaddr_t),
                           (const void *)b);
    return node;
}
\end{lstlisting}

\begin{lstlisting}
int nodeid_dump (uint8_t *b, const struct nodeID *s,
                 size_t max_write_size)
{
    return sockaddr_dump((void *)b, max_write_size, &s->addr);
}
\end{lstlisting}

\begin{lstlisting}
int node_ip(const struct nodeID *s, char *ip, int len)
{
    return sockaddr_strrep(&s->addr, ip, len) == NULL ? 0 : -1;
}
\end{lstlisting}

\begin{lstlisting}
int node_ip(const struct nodeID *s, char *ip, int len)
{
    return sockaddr_strrep(&s->addr, ip, len) == NULL ? 0 : -1;
}
\end{lstlisting}


\end{appendices}

\bibliographystyle{plain}
\begin{thebibliography}{99}

\bibitem{disi10-038} \textit{Luca Abeni, Csaba Kiraly, Alessandro Russo,
Marco Biazzini and Renato Lo Cigno},
June 2010, \textbf{GRAPES: a Generic Environment for P2P Streaming.}
{\em Technical report \#DISI-10-38}

\end{thebibliography}

\tableofcontents
\listoffigures
\listoftables

%\bibliography{rep-bib}

\end{document}
